% !TEX root = ../algorithms.tex

\section{Потоки в сетях. Алгоритм Форда--Фалкерсона. Теорема о максимальном потоке и минимальном разрезе}

\begin{Definition}{Транспортная сеть}{transport-network}
	Пусть дан ориентированный граф \(G=(V,E,c)\), где функция
	\[
	c\colon E \to \RR_{\ge 0}
	\]
	задаёт пропускные способности рёбер.
	
	Граф \(G\) называется \emph{транспортной сетью}, если в нём выделены две вершины:
	\[
	s \in V \quad \text{(исток)}, \qquad
	t \in V \quad \text{(сток)},
	\]
	причём
	\[
	\deg^{-}(s)=0, \qquad \deg^{+}(t)=0.
	\]
\end{Definition}

% ============================================================
%  Рис. 1 — Пример транспортной сети
% ============================================================
\begin{figure}[H]
	\centering
	\begin{tikzpicture}[arr]
		% --- вершины ---
		\node[nd source] (s) at (0,    0)   {$s$};
		\node[nd]        (a) at (2.8,  1.6) {};
		\node[nd]        (b) at (6,    1.6) {};
		\node[nd]        (c) at (2.8, -1.6) {};
		\node[nd]        (d) at (6,   -1.6) {};
		\node[nd source] (t) at (8.5,  0)   {$t$};
		
		\draw[arr] (s) -- node[above left,  lbl] {$2$}           (a);
		\draw[arr] (s) -- node[below left,  lbl] {$2$}           (c);
		\draw[arr] (a) -- node[above,       lbl] {$9$}           (b);
		\draw[arr] (a) to[bend left=18]
		node[right, pos=0.45, lbl] {$5$}       (c);
		\draw[arr] (c) to[bend left=18]
		node[left,  pos=0.45, lbl] {$4$}       (a);
		\draw[arr] (b) -- node[above right, lbl] {$1$}           (t);
		\draw[arr] (c) -- node[below,       lbl] {$3$}           (d);
		\draw[arr] (d) -- node[below right, lbl] {$3$}           (t);
	\end{tikzpicture}
	\caption{Пример транспортной сети. Числа у рёбер — пропускные способности $c(u,v)$.}
	\label{fig:network}
\end{figure}

\begin{Definition}{Поток}{}
	Потоком в транспортной сети называется функция $f\colon E \to \RR_{\ge 0}$, такая что:
	\begin{enumerate}[label=\arabic*)]
		\item \textbf{Ограничение пропускной способности:}
		\[
		0 \le f(u,v) \le c(u,v)
		\qquad \forall (u,v)\in E.
		\]
		\item \textbf{Закон сохранения потока:} для любой вершины $u \in V \setminus \{s,t\}$ выполнено
		\[
		\sum_{(v,u)\in E} f(v,u)
		=
		\sum_{(u,w)\in E} f(u,w).
		\]
	\end{enumerate}
\end{Definition}

% ============================================================
%  Рис. 2 — Пример потока в той же сети
%           Формат на рёбрах: пропускная способность | поток
% ============================================================
\begin{figure}[htb]
	\centering
	\begin{tikzpicture}[arr]
		\node[nd source] (s) at (0,    0)   {$s$};
		\node[nd]        (a) at (2.8,  1.6) {};
		\node[nd]        (b) at (6,    1.6) {};
		\node[nd]        (c) at (2.8, -1.6) {};
		\node[nd]        (d) at (6,   -1.6) {};
		\node[nd source] (t) at (8.5,  0)   {$t$};
		
		\draw[arr] (s) -- node[above left, lbl] {$2\mid\textcolor{blue}{2}$} (a);
		\draw[arr] (s) -- node[below left, lbl] {$2\mid\textcolor{blue}{2}$} (c);
		\draw[arr] (a) -- node[above, lbl] {$9\mid\textcolor{blue}{1}$} (b);
		\draw[arr] (a) to[bend left=18]
		node[right, pos=0.45, lbl] {$5\mid\textcolor{blue}{1}$} (c);
		\draw[arr] (c) to[bend left=18]
		node[left, pos=0.45, lbl] {$4\mid\textcolor{blue}{0}$} (a);
		\draw[arr] (b) -- node[above right, lbl] {$1\mid\textcolor{blue}{1}$} (t);
		\draw[arr] (c) -- node[below, lbl] {$3\mid\textcolor{blue}{3}$} (d);
		\draw[arr] (d) -- node[below right, lbl] {$3\mid\textcolor{blue}{3}$} (t);
		
		\node[lbl, align=left] at (9.6, 1.5) {%
			$c(e)\mid{\color{blue}f(e)}$\\[2pt]
			$\lvert f \rvert = 4$};
	\end{tikzpicture}
	\caption{Поток величины $4$ в транспортной сети из рис.~\ref{fig:network}.
		Рядом с каждым ребром указана пара
		\emph{пропускная способность\,|\,поток}.}
	\label{fig:flow}
\end{figure}

\begin{Definition}{Величиной потока }{}
	Величиной потока $f$ называется число
	\[
	|f|=\sum_{(s,v)\in E} f(s,v).
	\]
	Эквивалентно,
	\[
	|f|=\sum_{(u,t)\in E} f(u,t).
	\]
\end{Definition}

\begin{Proposition}{}{}
	Для любого потока $f$ выполнено
	\[
	\sum_{(s,v)\in E} f(s,v)=\sum_{(u,t)\in E} f(u,t).
	\]
\end{Proposition}

\begin{proof}
	Просуммируем закон сохранения потока по всем вершинам из множества $V\setminus\{s,t\}$. Все слагаемые, соответствующие внутренним рёбрам, сократятся попарно, и останется только суммарный поток, выходящий из истока, и суммарный поток, входящий в сток.
\end{proof}

\subsection{Задача о максимальном потоке}

Требуется найти поток максимальной величины.

\begin{Remark}{}{}
	Многие комбинаторные задачи допускают сведение к задаче о максимальном потоке.
\end{Remark}

\begin{Remark}{}{}
	Будем считать, что в графе нет кратных рёбер. Это не ограничивает общность: несколько параллельных рёбер можно заменить одним ребром с суммарной пропускной способностью.
\end{Remark}

% ============================================================
%  Рис. 3 — Сведение кратных рёбер
% ============================================================
\begin{figure}[htb]
	\centering
	\begin{tikzpicture}[arr, node distance = 2.8cm]
		% --- Левая часть ---
		\node[nd] (A1) at (0,0) {};
		\node[nd] (B1) at (3.2,0) {};
		
		% Два ребра A1 -> B1
		\draw[arr, bend left=45]  (A1) to node[above left, lbl] {$3$} (B1);
		\draw[arr, bend left=12]  (A1) to node[above, lbl] {$2$} (B1);
		
		% Одно ребро B1 -> A1
		\draw[arr, bend left=45]  (B1) to node[below right, lbl] {$1$} (A1);
		
		% --- Стрелка преобразования ---
		\node[font=\Large] at (4.8,0) {$\Longrightarrow$};
		
		% --- Правая часть ---
		\node[nd] (A2) at (6.5,0) {};
		\node[nd] (B2) at (9.7,0) {};
		
		% Суммарное ребро A2 -> B2
		\draw[arr, bend left=35] (A2) to node[above, lbl] {$5$} (B2);
		
		% Обратное ребро B2 -> A2
		\draw[arr, bend left=35] (B2) to node[below, lbl] {$1$} (A2);
	\end{tikzpicture}
	
	\label{fig:parallel}
\end{figure}
\begin{Proposition}{}{}
	Задача о максимальном потоке допускает формулировку в виде задачи линейного программирования:
	\[
	\max \sum_{(s,v)\in E} f(s,v)
	\]
	при ограничениях
	\[
	0 \le f(u,v) \le c(u,v)
	\qquad \forall (u,v)\in E,
	\]
	\[
	\sum_{(v,u)\in E} f(v,u)-\sum_{(u,w)\in E} f(u,w)=0
	\qquad \forall u\in V\setminus\{s,t\}.
	\]
\end{Proposition}

\begin{Remark}{}{}
	Множество допустимых потоков замкнуто и ограничено в $\mathbb{R}^{E}$, следовательно, компактно. Поэтому непрерывная функция
	\[
	(f(e))_{e\in E}\mapsto \sum_{(s,v)\in E} f(s,v)
	\]
	достигает на нем максимума по теореме Вейерштрасса.
\end{Remark}

\subsection{Идея построения максимального потока}

\begin{enumerate}[label=\arabic*)]
	\item Начинаем с нулевого потока $f\equiv 0$.
	\item На каждом шаге пытаемся увеличить поток, пропуская дополнительную порцию потока из $s$ в $t$ по некоторому допустимому пути.
\end{enumerate}

Однако возникает существенная трудность: иногда для увеличения общего потока необходимо не только добавлять поток по прямым рёбрам, но и частично \emph{отменять} ранее проведённый поток по некоторым рёбрам. Для формализации этого вводится остаточная сеть.

% ============================================================
%  Рис. 4 — Зачем нужны обратные рёбра (жадный сценарий)
% ============================================================
\begin{figure}[htb]
	\centering
	\begin{tikzpicture}[arr]
		
		% ---- Левая диаграмма: исходная сеть ----
		\begin{scope}
			\node[lbl, above] at (2, 2.2) {Исходная сеть};
			\node[nd source] (s1) at (0,   0)   {$s$};
			\node[nd]        (u1) at (2,   1.3) {$u$};
			\node[nd]        (v1) at (2,  -1.3) {$v$};
			\node[nd source] (t1) at (4,   0)   {$t$};
			
			\draw[arr] (s1) -- node[above left,  lbl] {$1\mid0$} (u1);
			\draw[arr] (s1) -- node[below left,  lbl] {$1\mid0$} (v1);
			\draw[arr] (u1) -- node[right,       lbl] {$1\mid0$} (v1);
			\draw[arr] (u1) -- node[above right, lbl] {$1\mid0$} (t1);
			\draw[arr] (v1) -- node[below right, lbl] {$1\mid0$} (t1);
		\end{scope}
		-
		\node at (5.5, 0) {$\longrightarrow$};
		
		\begin{scope}[xshift=7cm]
			\node[lbl, above, align=center] at (2, 2.2)
			{После пути $s\to u\to v\to t$};
			\node[nd source] (s2) at (0,   0)   {$s$};
			\node[nd]        (u2) at (2,   1.3) {$u$};
			\node[nd]        (v2) at (2,  -1.3) {$v$};
			\node[nd source] (t2) at (4,   0)   {$t$};
			
			\draw[arrG] (s2) -- node[above left,  lbl, black] {$\textcolor{green!55!black}{1\mid1}$} (u2);
			\draw[arrG] (u2) -- node[right,       lbl, black] {$\textcolor{green!55!black}{1\mid1}$} (v2);
			\draw[arrG] (v2) -- node[below right, lbl, black] {$\textcolor{green!55!black}{1\mid1}$} (t2);
			
			\draw[arr]  (s2) -- node[below left,  lbl] {$\textcolor{red!65!black}{1\mid0}$} (v2);
			\draw[arr]  (u2) -- node[above right, lbl] {$1\mid0$} (t2);
			
			% Пояснение
			\node[lbl, red!70!black, align=center] at (2, -2.5)
			{поток нельзя увеличить\\без отмены части $f$};
		\end{scope}
		
	\end{tikzpicture}
	\caption{Жадный выбор пути $s\to u\to v\to t$ насыщает ребро $v\to t$
		и блокирует путь $s\to v\to t$.
		Для получения оптимального потока величины~$2$ необходимо
		<<отменить>> единицу потока через ребро $u\to v$
		с помощью обратного ребра остаточной сети.}
	\label{fig:greedy}
\end{figure}

\subsection{Дополняющие пути и остаточная сеть}

\begin{Definition}{}{}
	Пусть $f$ --- поток в сети $G=(V,E,c)$. \emph{Остаточной сетью} называется ориентированный граф
	\[
	G_f=(V,E_f,c_f),
	\]
	который строится следующим образом. Для каждого ребра $(u,v)\in E$:
	\begin{enumerate}[label=\arabic*)]
		\item если $c(u,v)-f(u,v)>0,$ то в $E_f$ добавляется \emph{прямое} ребро $(u,v)$ остаточной пропускной способности
		\[
		c_f(u,v)=c(u,v)-f(u,v);
		\]
		\item если $f(u,v)>0,$ то в $E_f$ добавляется \emph{обратное} ребро $(v,u)$ остаточной пропускной способности
		\[
		c_f(v,u)=f(u,v).
		\]
	\end{enumerate}
\end{Definition}

% ============================================================
%  Рис. 5 — Построение остаточной сети G_f
% ============================================================
\begin{figure}[htb]
	\centering
	\begin{tikzpicture}[arr]
		
		% ===== Левая часть: сеть G с потоком =====
		\begin{scope}
			\node[lbl] at (2, 2.1) {$G$ (поток $|f|=1$)};
			
			\node[nd source] (s1) at (0,  0)    {$s$};
			\node[nd]        (u1) at (2,  1.3)  {$u$};
			\node[nd]        (v1) at (2, -1.3)  {$v$};
			\node[nd source] (t1) at (4,  0)    {$t$};
			
			% Насыщенные рёбра (зелёные)
			\draw[arrG] (s1) -- node[above left,  lbl, black] {$1\mid1$} (u1);
			\draw[arrG] (u1) -- node[right,       lbl, black] {$1\mid1$} (v1);
			\draw[arrG] (v1) -- node[below right, lbl, black] {$1\mid1$} (t1);
			
			% Ненасыщенные рёбра
			\draw[arr]  (s1) -- node[below left,  lbl] {$1\mid0$} (v1);
			\draw[arr]  (u1) -- node[above right, lbl] {$1\mid0$} (t1);
		\end{scope}
		
		% ===== Стрелка =====
		\node at (5.5, 0) {$\xrightarrow{\;G_f\;}$};
		
		% ===== Правая часть: остаточная сеть G_f =====
		\begin{scope}[xshift=7.2cm]
			\node[lbl] at (2, 2.1) {$G_f$};
			
			\node[nd source] (s2) at (0,  0)   {$s$};
			\node[nd]        (u2) at (2,  1.3) {$u$};
			\node[nd]        (v2) at (2, -1.3) {$v$};
			\node[nd source] (t2) at (4,  0)   {$t$};
			
			% Прямые рёбра (незасыщенные, синий пунктир)
			\draw[arrB] (s2) -- node[below left,  lbl, blue!65!black] {$1$} (v2);
			\draw[arrB] (u2) -- node[above right, lbl, blue!65!black] {$1$} (t2);
			
			% Обратные рёбра (красный пунктир)
			\draw[arrR] (u2) -- node[above left,  lbl, red!65!black]  {$1$} (s2);
			\draw[arrR] (v2) -- node[right,       lbl, red!65!black]  {$1$} (u2);
			\draw[arrR] (t2) -- node[below right, lbl, red!65!black]  {$1$} (v2);
			
			% Дополняющий путь s→v→u→t (жирный зелёный)
			\begin{scope}[on background layer]
				\draw[green!45!black, line width=4pt, opacity=0.35, rounded corners=4pt]
				(s2.center) -- (v2.center) -- (u2.center) -- (t2.center);
			\end{scope}
			
			% Легенда
			\node[lbl, align=left] at (6.1, 0) {%
				\textcolor{blue!65!black}{$\dashrightarrow$} прямое\\[3pt]
				\textcolor{red!65!black}{$\dashrightarrow$} обратное\\[3pt]
				\textcolor{green!55!black}{\rule{14pt}{3pt}} дополн. путь};
		\end{scope}
		
	\end{tikzpicture}
	\caption{Построение остаточной сети $G_f$.
		Синие пунктирные рёбра — прямые (можно добавить поток),
		красные пунктирные — обратные (можно отменить поток).
		Зелёным выделен дополняющий путь $s\to v\to u\to t$ в $G_f$,
		позволяющий увеличить поток до~$2$.}
	\label{fig:residual}
\end{figure}

\begin{Remark}{}{}
	В остаточной сети могут появляться кратные рёбра, если прямое и обратное направления возникают независимо. Это не влияет на корректность алгоритма.
\end{Remark}

\begin{Definition}{}{}
	Путь
	\[
	P=(v_0=s,v_1,\dots,v_k=t)
	\]
	в графе $G_f$ называется \emph{дополняющим путём} для потока $f$.
\end{Definition}

\begin{Definition}{}{}
	Пусть $P=(v_0=s,v_1,\dots,v_k=t)$ --- дополняющий путь в $G_f$. Его \emph{пропускной способностью} называется число
	\[
	\delta(P)=\min_{0\le i<k} c_f(v_i,v_{i+1}).
	\]
\end{Definition}

\begin{Proposition}{}{}
	Пусть $P=(v_0=s,v_1,\dots,v_k=t)$ --- дополняющий путь в $G_f$, и $\delta=\delta(P)$. Тогда можно построить новый поток $f'$ величины
	\[
	|f'|=|f|+\delta.
	\]
\end{Proposition}

\begin{proof}
	Для каждого ребра пути действуем следующим образом:
	\begin{enumerate}[label=\arabic*)]
		\item если ребро $(v_i,v_{i+1})$ является прямым ребром исходной сети, то полагаем
		\[
		f'(v_i,v_{i+1})=f(v_i,v_{i+1})+\delta;
		\]
		\item если ребро $(v_i,v_{i+1})$ является обратным, то есть в исходной сети имеется ребро $(v_{i+1},v_i)$, то полагаем
		\[
		f'(v_{i+1},v_i)=f(v_{i+1},v_i)-\delta.
		\]
	\end{enumerate}
	На всех остальных рёбрах поток не меняется.
	
	Поскольку $\delta \le c_f(v_i,v_{i+1})$ для каждого ребра пути, ограничения пропускной способности сохраняются. Закон сохранения потока также сохраняется: в каждой промежуточной вершине изменение входящего потока компенсируется изменением исходящего потока на ту же величину~$\delta$.
	
	Остаётся проверить, что величина потока увеличилась на $\delta$. По определению величины потока
	\[
	|f|=\sum_{(s,v)\in E} f(s,v).
	\]
	Первое ребро дополняющего пути, выходящее из $s$, является прямым ребром исходной сети, поэтому на этом ребре поток увеличивается на $\delta$. Другие изменения, связанные с вершиной $s$, отсутствуют. Следовательно,
	\[
	|f'|=\sum_{(s,v)\in E} f'(s,v)=\sum_{(s,v)\in E} f(s,v)+\delta=|f|+\delta.
	\]
\end{proof}

\subsection{Алгоритм Форда--Фалкерсона}

\begin{Statement}{}{}
	Алгоритм Форда--Фалкерсона строит поток итеративно следующим образом:
	\begin{enumerate}[label=\arabic*)]
		\item положить $f\equiv 0$;
		\item построить остаточную сеть $G_f$;
		\item найти в $G_f$ путь $P$ из $s$ в $t$;
		\item если такого пути нет, остановиться: текущий поток $f$ является максимальным;
		\item если путь $P$ найден, увеличить поток вдоль него на величину $\delta(P)$, получив новый поток $f$;
		\item перестроить остаточную сеть $G_f$ для нового потока и перейти к шагу 3.
	\end{enumerate}
\end{Statement}

\begin{Remark}{}{}
	На шаге поиска пути можно использовать как DFS, так и BFS. От выбора способа поиска существенно зависит итоговая оценка времени работы.
\end{Remark}

\subsection{\texorpdfstring{$s$--$t$}{s-t}-разрез}

\begin{Definition}{}{}
	\emph{$s$--$t$-разрезом} называется разбиение множества вершин $V=S \sqcup T,$ 
	$\quad s\in S, t\in T.$
\end{Definition}

\begin{Definition}{}{}
	\emph{Величиной разреза} $(S,T)$ называется сумма пропускных способностей всех рёбер, идущих из $S$ в $T$:
	\[
	c(S,T)=\sum_{\substack{(u,v)\in E\\ u\in S,\ v\in T}} c(u,v).
	\]
\end{Definition}

\begin{figure}[htb]
	\centering
	\begin{tikzpicture}[arr, scale=1, every node/.style={transform shape}]
		
		% ===== Вершины =====
		\node[nd source] (s)  at (0,0) {$s$};
		\node[nd]        (u1) at (2.2,1.35) {};
		\node[nd]        (u2) at (2.2,-1.35) {};
		
		\node[nd]        (v1) at (5.2,1.35) {};
		\node[nd]        (v2) at (5.2,-1.35) {};
		\node[nd source] (t)  at (7.4,0) {$t$};
		
		% ===== Фоны множеств S и T =====
		\begin{scope}[on background layer]
			\fill[blue!8, draw=blue!55, thick, rounded corners=18pt]
			(-0.55,-2.05) rectangle (2.85,2.05);
			\fill[orange!8, draw=orange!65, thick, rounded corners=18pt]
			(4.55,-2.05) rectangle (7.95,2.05);
		\end{scope}
		
		\node[blue!70!black, font=\small\bfseries]   at (1.15, 1.85) {$S$};
		\node[orange!70!black, font=\small\bfseries] at (6.25, 1.85) {$T$};
		
		% ===== Рёбра внутри S =====
		\draw[arr, gray!65] (s) -- (u1);
		\draw[arr, gray!65] (s) -- (u2);
		\draw[arr, gray!65] (u1) to[bend right=18] (u2);
		
		% ===== Рёбра внутри T =====
		\draw[arr, gray!65] (v1) -- (t);
		\draw[arr, gray!65] (v2) -- (t);
		\draw[arr, gray!65] (v1) to[bend right=18] (v2);
		
		% ===== Рёбра разреза S -> T =====
		\draw[arrC] (u1) -- node[above, lbl, pos=0.52, yshift=2pt, black] {$c_1$} (v1);
		\draw[arrC] (u1) -- node[above, lbl, pos=0.32, xshift=3pt, yshift=-2pt, black] {$c_2$} (v2);
		\draw[arrC] (u2) -- node[below, lbl, pos=0.50, yshift=-2pt, black] {$c_3$} (v2);
		
		% ===== Обратное ребро T -> S (не учитывается) =====
		\draw[arr, gray!50] (v1) to[bend right=36]
		node[lbl, gray!70, pos=0.42, right=2pt] {$c'$} (u2);
		
		% ===== Пунктирная линия разреза =====
		\draw[thick, densely dashed, gray!60] (3.7,2.55) -- (3.7,-2.75);
		\node[font=\footnotesize, gray!70] at (3.7,-3.02) {разрез};
		
		% ===== Формула =====
		\node[lbl, align=center] at (3.7,-4.15) {%
			$c(S,T)=c_1+c_2+c_3$\\[2pt]
			(ребро $c'\colon T\to S$ не учитывается)};
		
	\end{tikzpicture}
	\caption{$s$--$t$-разрез $(S,T)$ и его величина $c(S,T)$. 
		Красные рёбра идут из $S$ в $T$ и входят в сумму, 
		а серое ребро из $T$ в $S$ не учитывается.}
	\label{fig:cut}
\end{figure}

\begin{Proposition}{}{}
	Для любого потока $f$ и любого $s$--$t$-разреза $(S,T)$ выполнено неравенство
	\[
	|f|\le c(S,T).
	\]
\end{Proposition}

\begin{proof}
	Рассмотрим сумму
	\[
	\sum_{u\in S}
	\left(
	\sum_{(v,u)\in E} f(v,u)
	-
	\sum_{(u,w)\in E} f(u,w)
	\right).
	\]
	Для всех вершин из $S\setminus\{s\}$ эта величина равна нулю по закону сохранения потока, а для вершины $s$ она равна $-|f|$. После сокращения внутренних рёбер получаем
	\[
	|f|=
	\sum_{\substack{(u,v)\in E\\ u\in S,\ v\in T}} f(u,v)
	-
	\sum_{\substack{(u,v)\in E\\ u\in T,\ v\in S}} f(u,v).
	\]
	Следовательно,
	\[
	|f|
	\le
	\sum_{\substack{(u,v)\in E\\ u\in S,\ v\in T}} f(u,v)
	\le
	\sum_{\substack{(u,v)\in E\\ u\in S,\ v\in T}} c(u,v)
	=
	c(S,T).
	\]
\end{proof}

\subsection{Теорема о максимальном потоке и минимальном разрезе}

\begin{Theorem}{}{}
	Для потока $f$ в транспортной сети эквивалентны следующие условия:
	\begin{enumerate}[label=\arabic*)]
		\item $f$ является максимальным потоком;
		\item в остаточной сети $G_f$ нет дополняющего пути из $s$ в $t$;
		\item существует $s$--$t$-разрез $(S,T)$ такой, что
		\[
		|f|=c(S,T).
		\]
	\end{enumerate}
\end{Theorem}

\begin{proof}
	Докажем циклически
	\[
	1)\Rightarrow 2)\Rightarrow 3)\Rightarrow 1).
	\]
	
	\textbf{$1)\Rightarrow 2)$.}
	Если бы в $G_f$ существовал дополняющий путь, то по предыдущему предложению поток можно было бы увеличить. Это противоречит максимальности $f$.
	
	\textbf{$2)\Rightarrow 3)$.}
	Пусть $S$ --- множество вершин, достижимых из $s$ в остаточной сети $G_f$, а
	\[
	T=V\setminus S.
	\]
	Тогда $s\in S$, а $t\notin S$, поскольку по предположению пути из $s$ в $t$ в $G_f$ нет.
	
	Покажем, что
	\[
	|f|=c(S,T).
	\]
	Из доказательства оценки для разреза имеем
	\[
	|f|=
	\sum_{\substack{(u,v)\in E\\ u\in S,\ v\in T}} f(u,v)
	-
	\sum_{\substack{(u,v)\in E\\ u\in T,\ v\in S}} f(u,v).
	\]
	Теперь используем определение множества $S$.
	
	Если $(u,v)\in E$, $u\in S$, $v\in T$, то в остаточной сети не может быть прямого ребра $(u,v)$, иначе вершина $v$ была бы достижима из $s$. Значит,
	\[
	c(u,v)-f(u,v)=0,
	\qquad\text{то есть}\qquad
	f(u,v)=c(u,v).
	\]
	
	Если $(u,v)\in E$, $u\in T$, $v\in S$, то в остаточной сети не может быть обратного ребра $(v,u)$, иначе вершина $u$ была бы достижима из $s$. Значит,
	\[
	f(u,v)=0.
	\]
	
	Следовательно,
	\[
	|f|=
	\sum_{\substack{(u,v)\in E\\ u\in S,\ v\in T}} c(u,v)
	=
	c(S,T).
	\]
	
	\textbf{$3)\Rightarrow 1)$.}
	Для любого потока $f'$ и любого разреза $(S,T)$ верно
	\[
	|f'|\le c(S,T).
	\]
	Если для некоторого потока $f$ выполнено равенство
	\[
	|f|=c(S,T),
	\]
	то никакой поток не может иметь величину больше $|f|$. Следовательно, $f$ максимален.
\end{proof}

\begin{Proposition}{}{}
	Максимальная величина потока равна минимальной величине $s$--$t$-разреза.
\end{Proposition}

\begin{proof}
	Теорема утверждает существование потока и разреза, для которых достигается равенство
	\[
	|f|=c(S,T).
	\]
	С другой стороны, для любого потока и любого разреза всегда выполнено неравенство
	\[
	|f|\le c(S,T).
	\]
	Отсюда и следует утверждение.
\end{proof}

\subsection{Оценки времени работы}

\begin{Proposition}{}{}
	Одна итерация алгоритма Форда--Фалкерсона работает за время
	\[
	O(|E|),
	\]
	если поиск дополняющего пути выполняется обходом графа в глубину или в ширину.
\end{Proposition}

\begin{proof}
	Построение остаточной сети и поиск пути в ней занимают линейное по числу рёбер время.
\end{proof}

\begin{Proposition}{}{}
	Если все пропускные способности целочисленны, то алгоритм Форда--Фалкерсона завершится не более чем за
	\[
	|f_{\max}|
	\]
	итераций, где $f_{\max}$ --- максимальный поток.
\end{Proposition}

\begin{proof}
	При целочисленных пропускных способностях каждая остаточная пропускная способность также целочисленна. Следовательно, на каждой итерации
	\[
	\delta(P)\ge 1,
	\]
	и величина потока увеличивается хотя бы на единицу.
\end{proof}

\begin{Corollary}{}{}
	В случае целочисленных пропускных способностей время работы алгоритма Форда--Фалкерсона оценивается как
	\[
	O(|E|\cdot |f_{\max}|).
	\]
\end{Corollary}

\begin{Remark}{}{}
	Для произвольных вещественных пропускных способностей алгоритм Форда--Фалкерсона при неудачном выборе дополняющих путей может работать очень долго; в общем случае возможны патологические сценарии.
\end{Remark}

\begin{Statement}{}{}
	Если на каждом шаге искать кратчайший по числу рёбер дополняющий путь с помощью BFS, то получается алгоритм Эдмондса--Карпа со временем работы
	\[
	O(|V|\cdot |E|^2).
	\]
\end{Statement}

\begin{Remark}{}{}
	Существуют и более быстрые алгоритмы решения задачи о максимальном потоке, в частности алгоритм Диница и алгоритм Орлина.
\end{Remark}
